\documentclass[10pt,twoside]{article}
\usepackage[utf8]{inputenc}
\usepackage{amsmath}
\usepackage{amsfonts}
\usepackage{amssymb}
\usepackage[spanish,es-noshorthands]{babel}
\usepackage[T1]{fontenc}
\usepackage{lmodern}
\usepackage{graphicx,hyperref}
\usepackage{tikz,pgf}
\usepackage{multicol}
\usepackage{subfig}
\usepackage[papersize={6.5in,8.5in},width=5.5in,height=7in]{geometry}
\usepackage{fancyhdr}
\pagestyle{fancy}
\fancyhead[LE]{\includegraphics[height=12pt]{Images/logo-colegio.png} Matemáticas $11^{\circ}$}
\fancyhead[RE]{}
\fancyhead[RO]{\textit{Germ\'an Avenda\~no Ram\'irez, Lic. U.D., M.Sc. U.N.}}
\fancyhead[LO]{}

\author{Germ\'an Avenda\~no Ram\'irez, Lic. U.D., M.Sc. U.N.}
\title{\begin{minipage}{.2\textwidth}
\includegraphics[height=1.75cm]{Images/logo-colegio.png}\end{minipage}
\begin{minipage}{.55\textwidth}
\begin{center}
Taller de Nivelación 2014\\
Matemáticas $11^{\circ}$
\end{center}
\end{minipage}\hfill
\begin{minipage}{.2\textwidth}
\includegraphics[height=1.75cm]{Images/logo-sed.png} 
\end{minipage}}
\date{}
\begin{document}
\maketitle
Nombre: \hrulefill Curso: \underline{\hspace*{44pt}} Fecha: \underline{\hspace*{2.5cm}}
\section*{N\'{u}meros reales}
\begin{enumerate}
\item 
\begin{enumerate}
\item Grafique el intervalo $(-5,3)$ y $(2,\infty)$ en la recta real
\item Exprese las desigualdades $x\leq 3$ y $-1\leq x<4$ en notación de intervalos
\item Encuentre la distancia entre $-7$ y $9$ sobre la recta real
\end{enumerate}
\item Evalúe cada expresión
\begin{enumerate}
\begin{multicols}{5}
\item $(-3)^{4}$ \item $-3^{4}$ \item $\dfrac{5^{23}}{5^{24}}$ \item $\left(\dfrac{3}{3}\right)^{-2}$ \item $16^{-3/4}$
\end{multicols}
\end{enumerate}
\item Escriba cada n\'{u}mero en notaci\'{o}n cient\'{i}fica
\begin{enumerate}
\begin{multicols}{2}
\item 186\,000'000\,000 \item 0.0000003965
\end{multicols}
\end{enumerate}
\item Simplifique cada expresión. Escriba su respuesta final sin exponentes negativos
\begin{enumerate}
\begin{multicols}{3}
\item $\sqrt{200}-\sqrt{32}$ \item $(3a^{3}b^{3})(4ab^{2})^{2}$ \item $\left(\dfrac{3x^{3/2}y^{3}}{x^{2}y^{-1/2}}\right)^{-2}$
\item $\dfrac{x^{2}+3x+2}{x^{2}-x-2}$
\item $\dfrac{x^{2}}{x^{2}-4}-\dfrac{x+1}{x+2}$
\item $\dfrac{\frac{y}{x}-\frac{x}{y}}{\frac{1}{y}-\frac{1}{x}}$
\end{multicols}
\end{enumerate}
\item Racionalice el denominador y simplifique: $\dfrac{\sqrt{10}}{\sqrt{5}-2}$
\item Realice las operaciones indicadas y simplifique:
\begin{enumerate}
\begin{multicols}{3}
\item $3(x+6)+4(2x-5)$
\item $(x+3)(4x-5)$
\item $(\sqrt{a}+\sqrt{b})(\sqrt{a}-\sqrt{b})$
\item $(2x+3)^{2}$
\item $(x+2)^{3}$
\end{multicols}
\end{enumerate}
\item Factorice completamente cada expresión
\begin{enumerate}
\begin{multicols}{3}
\item $4x^{2}-25$
\item $2x^{2}+5x-12$
\item $x^{3}-3x^{2}-4x+12$
\item $x^{4}+27x$
\item $3x^{3/2}-9x^{1/2}+6x^{-1/2}$
\item $x^{3}y-4xy$
\end{multicols}
\end{enumerate}
\item Encuentre las soluciones reales:
\begin{enumerate}
\begin{multicols}{3}
\item $x+5=14-\frac{1}{2}x$
\item $\dfrac{2x}{x+1}=\dfrac{2x-1}{x}$
\item $x^{2}-x-12=0$
\item $2x^{2}+4x+1=0$
\item $\sqrt{3-\sqrt{x+5}}=2$
\item $x^{4}-3x^{2}+2=0$
\item $3|x-4|=10$
\end{multicols}
\end{enumerate}
\item Mary condujo de Bogotá a Melgar a una rapidez promedio de 80 km/h. De regreso, ella condujo en promedio a 70 km/h. El tiempo total de viaje fue de $4\frac{2}{3}$ de hora. Encuentre la distancia entre las dos ciudades.
\item Una lote rectangular tiene 70 m más de largo que de ancho y su diagonal mide 130 m. Encuentre las dimensiones del lote.
\item Solucione cada inecuación. Escriba la respuesta usando la notación de intervalos y dibuje la solución en la recta real.
\begin{enumerate}
\begin{multicols}{2}
\item $-4<5-3x\leq 17$
\item $x(x-1)(x+2)>0$
\item $|x-4|<3$
\item $\dfrac{2x-3}{x+1}\leq 1$
\end{multicols}
\end{enumerate}
\item Una botella de medicina debe ser guardada a una temperatura entre $5^{\circ}$C y 10$^{\circ}$C. Qué rango correspondería si se toma la escala Fahrenheit? (Recuerde que la temperatura en Fahrenheit (F) y Celsius (C) satisface la relación $C=\frac{5}{9}(F-32)$
\end{enumerate}
\end{document}
